\chapter{Deep noise suppression}

\section{Metrics}
Standard practice measures speech quality with intrusive and non-intrusive metrics. For intrusive \ac{SE} metrics, we measure PESQ~\citep{rixPerceptualEvaluationSpeech2001} for predicting speech quality, ESTOI~\citep{jensen2016algorithm} as a measure of speech intelligibility and \ac{SI-SDR}~\citep{leroux2018sdr} measured in dB. We also measure non-intrusive metrics that predict quality from the predicted clean speech alone. Firstly, we compute the common metric DNSMOS~\citep{reddy2021dnsmos},\footnote{https://github.com/microsoft/DNS-Challenge/tree/master/DNSMOS} which employs a neural network trained on human ratings (\ac{MOS}).  Secondly, we use WhiSQA, a non-intrusive \ac{MOS} prediction network shown to correlate well with human judgment \citep{closeHallucinationPerceptualMetricDriven2024,close2025whisqanonintrusivespeechquality}.\footnote{https://github.com/leto19/WhiSQA} All of the above metrics score higher for better quality speech.

\section{Data}
We train and test all experiments on the \ac{VB-DMD} dataset~\citep{valentini2016investigating}, a common benchmark for \ac{SE} containing clean speech recordings from 28 speakers with added background noise, e.g. caf\'e, traffic. We use speakers p226 and p287 for validation. Non-intrusive evaluation of the clean speech yields 3.53 DNSMOS and 4.53 WhiSQA. 

\section{Model}
Following \cite{jukic2024schr}, we train all models until validation SI-SDR converges, then choose the checkpoint with the best validation PESQ. Unless stated, we run \ac{ODE} samplers for 50 steps, with batch size 8 and the same STFT settings as \cite{richterInvestigatingTrainingObjectives2025}.
The neural estimator $F_\theta$ employs the NCSN++ architecture \citep{song2021sde} using the same parameterisation described in \cite{richter2023speech}. All experiments use the time-domain auxiliary loss \citep{jukic2024schr}.


The neural estimator $F_\theta$ employs the NCSN++ architecture \citep{song2021sde} using the same parameterisation described in \cite{richter2023speech}. All experiments use the time-domain auxiliary loss \citep{jukic2024schr}.

\begin{table}[htbp]
\centering
\resizebox{\linewidth}{!}{%
\begin{tabular}{cllll|ccccc}
        \toprule
               Path& Loss&Inference& k&$c$&  PESQ &  ESTOI &  SI-SDR&   DNSMOS&WhiSQA \\
    \midrule
     Noisy& -&-&-&-& 1.97& 0.79& 8.4&  3.05&3.11\\
             \midrule
    SB-VE& DP&ODE& 2.6&0.4& 2.92& 0.87& 19.3& 3.56 & 4.46\\
     SB-VE&  DP&ODE& 0.99& 0.375& 2.92 & 0.88 & 19.5& 3.56 &4.47 \\
             \midrule
    SB-SV& DP&ODE& 2.6&0.15& 2.98 & 0.88 & 19.4 & 3.58   & \textbf{4.51   }\\
    SB-SV& DP&ODE& 0.99&0.1& 2.86 & 0.88 & 19.5 & 3.58   & 4.47   \\
    ICFM& DP&ODE& -&0.1& 2.98 & 0.88 & 20.1 & 3.59  & 4.49\\
    ICFM& FM&ODE& -&0.1& 2.91 & 0.88 & 20.3 & \textbf{3.60} & 4.50 \\
    \midrule
             SB-VE   & DP&DDP& 2.6&0.4&  2.92&  0.87&  19.4&   3.55&4.45\\
             SB-SV& DP&DDP& 2.6&0.15& 2.98& 0.88& 19.9&  3.58&4.50\\
             SB-SV& DP&DDP& 0.99&0.1& 2.99& 0.88& 20.0&  3.57&4.50\\
             ICFM& DP&DDP& -&0.1&  \textbf{3.05}&  0.88&  20.2&   3.58&\textbf{4.51}\\
            ICFM& FM&DDP& -&0.1&  3.00&  0.88&  \textbf{20.4}&   3.59&\textbf{4.51}\\
            \bottomrule
        \end{tabular}
}
\caption{Mean speech quality metrics on \ac{VB-DMD} of our SB-SV and ICFM with FM loss and \ac{DDP} inference over SB-VE \citep{jukic2024schr} baseline. $k=0.99$ induces straightness and $c$ scales variance.}
\label{tab:results}
\end{table}
Our results are displayed in \tableref{tab:results}. Our proposed straighter paths \ac{SB-SV} (\sectionref{sec:sbsv}) and \ac{ICFM} (\sectionref{sec:icfm}) suggest improved speech quality across all metrics over the curved SB-VE (\sectionref{sec:schr}). Interestingly, there is no apparent benefit of using SB-SV or SB-VE with a more linear path ($k=0.99$). In fact, PESQ and WhiSQA decrease when using SB-SV with $k=0.99$.  However, compared to SB-SV, the results show that using an exact linear gradient with static variance with ICFM produces higher quality samples. Using ICFM with the FM loss has a marginal improvement over DP. This difference in training objectives suggests that direct gradient estimation is more suitable for ICFM. Together, these findings support the idea that straighter paths are more suitable for flow-based SE, specifically by introducing time-independent variance.
%
\figureref{fig:N} shows how the number of ODE steps affects the performance of various model types from \tableref{tab:results}. Although SB-VE performs better at 1 ODE step, ICFM requires 20 steps to outperform SB-VE. SB-SV has a similar trend to ICFM, suggesting that static variance reduces performance when using fewer ODE steps. We speculate that models with static variance might perform worse with one-step ODEs because they don't exactly interpolate the data~\eqref{eq:icfm_mean_variance}, unlike SB-VE \eqref{eq:sb_mean_variance}.
%
On average, the samples of \ac{DDP} are either comparable to or of improved quality over those predicted with 50 ODE steps. Further, possible ODE errors in SB-SV $k=0.99$ are circumvented by DDP. Our proposed ICFM with FM loss reports the highest PESQ and SI-SDR. The results suggest that, although trained for ODE solvers, flow-based models have prominent predictive properties. Another reason ICFM performs well here could be attributed to variance at the boundaries, alleviating potential overfitting caused by exact interpolation.

\section{Conclusion}
 This chapter views the time-independence of path gradient and variance as an analogue for straightness. We assessed the impact of probability path straightness on flow-based model performance for SE. By comparing SB-VE with SB-SV, we observed greater improvement with time-independent variance over gradient, but overall found that speech quality metrics were greater improved by using ICFM, which fixes both gradient and variance. However, fixing variance degraded ODE solver performance, but this can be circumvented by directly predicting the data at inference.